\section{L'information qui consomme~: convergence numérique et substrat énergétique (2007--)}\label{sec:contemporain}
% =============================================================================

\subsection{Le \emph{turning point} et l'invention du surplus comportemental (2001--2008)}\label{subsec:turning_point}

\dimmark{D5/D6}{+}%
La section précédente s'est achevée sur une tension irrésolue~: la phase \emph{frenzy} de la fin des années 1990 avait déployé une infrastructure et capté une base d'utilisateurs sans qu'un modèle commercial stable ne stabilise les flux financiers. \textcite{goldhaber_attention_1997} avait nommé le problème sans le résoudre. La résolution intervient entre 2001 et 2003 dans une firme alors périphérique à la bulle qui venait d'éclater~: Google découvre que les traces de comportement des utilisateurs~--- clics, requêtes, durées de session, abandons~--- ne sont pas seulement des données nécessaires à l'amélioration de son moteur de recherche, mais des actifs prédictifs valorisables sur un marché publicitaire qui se constitue par leur agrégation. La traduction opérationnelle est \emph{AdWords} (lancé en octobre 2000, repensé en 2002) puis \emph{AdSense} (2003)~: un mécanisme d'enchères sur l'attention des utilisateurs ciblés par leur historique comportemental. \textcite{zuboff_surveillance_2019} a proposé de cette séquence le terme de \emph{surplus comportemental}~: une catégorie d'actifs informationnels dérivés du comportement humain et capturables sans contrepartie marchande directe, dont l'agrégation forme la matière première d'un nouveau type de marché (\emph{behavioural futures markets}).

\dimmark{D5/D10}{+}%
L'année 2007--2008 marque la généralisation du mécanisme à l'échelle de l'écosystème. Facebook, fondé en 2004 et longtemps réticent à déployer la publicité par crainte de cannibaliser sa croissance utilisateur, adopte en 2007 un \emph{Ad Platform} comparable à celui de Google, puis publie en 2008 \emph{Facebook Beacon} qui tente d'étendre la capture aux comportements hors plateforme~--- une extension trop visible qui provoque un recul public mais inscrit le principe~\parencite{west_data_2019}. La même année, le lancement de l'iPhone en juin 2007 puis de l'App Store en juillet 2008 transforme l'équipement utilisateur en plateforme permanente de capture comportementale~: chaque application installée devient un capteur de données, chaque séquence d'interaction une enregistrement valorisable. \textcite{johnstone_deep_2026} identifient cette séquence 2001--2008 comme le \emph{turning point} de la trajectoire numérique au sens de \textcite{perez_technological_2002}, par lequel la phase \emph{frenzy} bascule vers la phase \emph{synergy}~: ce qui était dispersion spéculative devient configuration stable, ce qui était promesse devient infrastructure. La nuance que la présente analyse veut introduire concerne le contenu même de la stabilisation~: la résolution \emph{commerciale} de la tension fin 1990 (l'invention du modèle économique manquant) est aussi le moment où le couplage \emph{énergétique} cesse d'être ponctuel pour devenir continu à l'échelle planétaire, puisque la capture comportementale exige des serveurs en fonctionnement permanent. Les sections suivantes documentent cette continuité.


\subsection{L'informatique en nuage et le centre de données à très grande échelle}\label{subsec:cloud}

\dimmark{D2/D3}{+}%
Amazon Web Services ouvre ses premiers services en 2006, Google Compute Engine suit en 2012, Microsoft Azure s'impose dans les années suivantes~: en une décennie, le calcul cesse d'être un service local que l'entreprise produit pour elle-même et devient un flux acheté à des opérateurs d'infrastructure spécialisés. Le mouvement est structurellement analogue à celui qu'Edison et Insull avaient accompli entre 1882 et 1910~: remplacer des générateurs locaux dispersifs par une centrale mutualisée exploitant les économies d'échelle. L'entreprise qui maintenait sa propre salle informatique comme Edison avait ses lampes à arc propres cède la place à l'abonné du réseau~; le capital informatique fixe se transforme en coût variable d'exploitation. Le rendement d'échelle est ici triple~: le taux d'utilisation des serveurs monte de dix à soixante-dix pour cent par mutualisation des pic de charge~; le coût d'achat des composants baisse par l'effet du volume~; le coût d'opération (refroidissement, alimentation, maintenance) est optimisé à une échelle qu'aucune entreprise individuelle ne peut atteindre.

\dimmark{D8/0-TER}{+}%
La puissance électrique installée par ces centres de données devient un indicateur économique significatif. La consommation électrique mondiale des centres de données atteignait environ 415~TWh en 2024~--- soit 1,5~\% de la demande électrique mondiale~--- et l'Agence internationale de l'énergie projette une quasi-multiplication par deux d'ici 2030, à environ 945~TWh \parencite{iea_energy_ai_2025}. Pour situer l'ordre de grandeur~: c'est l'équivalent de la consommation électrique annuelle du Japon ajoutée en six ans. Dans certains états américains, la concentration géographique de l'infrastructure rend le coût sectoriel visible à l'échelle régionale~: en Virginie du Nord, les centres de données représentent plus de vingt-six pour cent de la consommation électrique régionale~; les files d'attente de raccordement au réseau de transport à haute tension atteignent sept ans pour les projets les plus importants.

\dimmark{D10}{+}%
Cette concentration crée un effet de dépendance de sentier du même type que le \emph{technological momentum} que \textcite{hughes_networks_1983} avait identifié pour le système électrique d'Insull~: des investissements massifs, des standards d'interconnexion, des institutions de régulation et des contrats d'approvisionnement énergétique s'accumulent autour d'une architecture technique et résistent collectivement au remplacement. Les quatre plus grands opérateurs~--- Amazon, Google, Microsoft, Meta~--- ont multiplié par trois leur consommation électrique entre 2018 et 2023 (de trente-cinq à plus de cent dix~TWh) et ont engagé plus de trois cent vingt milliards de dollars de dépenses d'investissement pour 2025 seul, un rythme qui n'a aucun précédent dans l'histoire du secteur informationnel.

\dimmark{0-TER}{+}%
L'empreinte agrégée du secteur déborde le seul poste électrique des centres de données. Sur la consommation énergétique opérationnelle, \textcite{bieser_review_2023} synthétisent les estimations disponibles et situent la contribution du secteur ICT aux émissions mondiales de gaz à effet de serre entre un et demi et quatre pour cent en~2023~; l'ampleur de l'intervalle, plus que la valeur centrale, est elle-même un résultat~: aucune méthodologie partagée ne stabilise le périmètre du secteur ni l'allocation de la consommation électrique entre équipements terminaux, réseaux et centres de données. Sur les flux de matière, \textcite{forti_global_2020} documentent que les déchets d'équipements électriques et électroniques (DEEE) atteignaient 53,6 millions de tonnes en 2019 et sont projetés à 74,7 millions de tonnes en 2030, dont les terminaux informatiques et de télécommunication constituent une part croissante. Sur l'angle climatique, \textcite{mazzucato_ai_2024} situent les émissions cumulées du \emph{cloud carbon} au même ordre de grandeur que celles du transport aérien commercial mondial, ordre de grandeur qui n'est pas une mesure consolidée mais une comparaison heuristique destinée à fixer l'échelle. Ces trois séries ne convergent pas vers un compte unifié~; elles convergent vers un constat~: l'empreinte matérielle du système informationnel contemporain est documentable, son ampleur est significative à l'échelle des grandes industries lourdes, et l'absence d'un cadre statistique partagé pour la mesurer est elle-même un fait analytique sur lequel le chapitre~\ref{ch:biophysique} aura à revenir.


\subsection{La photolithographie et le coût de la miniaturisation}\label{subsec:lithographie}

\dimmark{D4/D10}{+}%
Les sections précédentes ont documenté la consommation d'énergie du système informationnel du côté de l'usage~: centres de données, réseaux, terminaux. Elles n'ont pas examiné le substrat de \emph{production} qui rend possible l'existence même de ces équipements. Or la fabrication des semi-conducteurs, et en particulier la photolithographie qui en constitue l'étape la plus critique, est elle-même le lieu d'une intensification énergétique et matérielle considérable dont l'histoire éclaire la configuration actuelle du marché des puces.

La loi de Moore, formulée en 1965 comme observation empirique sur le doublement du nombre de transistors par circuit intégré, est devenue au cours des décennies suivantes une institution industrielle au sens fort du terme. \textcite{mody_long_2017} a montré que cette régularité n'est pas le résultat passif d'une trajectoire technique~; c'est une prophétie auto-réalisatrice, coordonnée par l'\emph{International Technology Roadmap for Semiconductors} (ITRS, 1991--2016, puis IRDS), autour de laquelle l'ensemble de l'industrie aligne ses cycles d'investissement, ses calendriers de R\&D et ses décisions d'équipement. La loi de Moore n'est pas une loi de la nature~; c'est un dispositif de coordination industrielle dont la force repose sur le fait que chaque acteur a intérêt à ce que les autres y croient.

\dimmark{0-TER}{*}%
La contrainte physique fondamentale de la miniaturisation est la résolution de la gravure optique. Les motifs du circuit sont projetés par un faisceau lumineux à travers un masque sur la surface du wafer de silicium~; la finesse des motifs gravables est bornée par la longueur d'onde de la lumière utilisée (limite de diffraction). Chaque génération de gravure a donc exigé de descendre en longueur d'onde~: lumière visible dans les années 1970, ultraviolet profond (DUV) à 248~nm puis 193~nm dans les années 1990 et 2000, et finalement ultraviolet extrême (EUV) à 13,5~nm à partir de la fin des années 2010. Le passage au DUV à immersion (193~nm avec lentille immergée dans l'eau) a permis de prolonger la trajectoire jusqu'au n{\oe}ud de 7~nm, au prix de techniques de \emph{multi-patterning} (gravure en plusieurs passes) qui multipliaient le nombre d'étapes de fabrication et donc le coût et la consommation d'énergie par couche gravée. Lorsque le multi-patterning est devenu économiquement prohibitif, seul le passage à l'EUV pouvait prolonger la trajectoire de miniaturisation.

\dimmark{D4/D8}{*}%
Ce passage a exigé trente ans de développement (1994--2024) et des investissements cumulés de l'ordre de plusieurs dizaines de milliards de dollars. L'industrialisation de l'EUV a été lancée en 1994 par une coalition de fabricants de semi-conducteurs, parmi lesquels Intel, TSMC et Samsung, associés à ASML, entreprise néerlandaise issue en 1984 d'une coentreprise entre ASM International et Philips. La source lumineuse EUV repose sur un mécanisme sans équivalent dans les générations précédentes~: un laser CO$_2$ de haute puissance vaporise des gouttelettes d'étain en plasma, qui émet un rayonnement à 13,5~nm~; ce rayonnement est collecté et focalisé par un système de miroirs multicouches (fabriqués par Carl Zeiss SMT en Allemagne) dont la surface doit être polie à une précision inférieure à l'atome. Le rendement de conversion énergétique de ce processus est de l'ordre de quatre à six pour cent~: plus de quatre-vingt-dix pour cent de l'énergie injectée est dissipée en chaleur\footnote{Les estimations de rendement varient selon les sources~; le chiffre de quatre à six pour cent correspond aux données publiées par les analystes de l'industrie pour les systèmes NXE~3400 en production vers 2020--2023. La génération suivante (High-NA, EXE~5000) utilise une optique à ouverture numérique de 0,55 au lieu de 0,33, ce qui améliore la résolution mais ne modifie pas fondamentalement le rendement de la source.}. Le refroidissement de la machine et de son environnement représente à lui seul jusqu'à trente pour cent de la consommation totale de l'installation.

\dimmark{D4}{*}%
Le résultat industriel de cette trajectoire est un monopole mondial sans équivalent dans l'histoire des technologies de production. ASML est aujourd'hui le seul fabricant de machines de lithographie EUV~; ses concurrents japonais Nikon et Canon, qui partageaient le marché DUV dans les années 1990, n'ont pas accompli la transition. La machine High-NA EUV (EXE~5200, livrée à partir de 2024) contient plus de cent mille composants, pèse l'équivalent de deux avions de ligne, et coûte environ 350~millions d'euros l'unité. Ses trois principaux clients~--- TSMC, Samsung et Intel~--- sont aussi les entreprises qui ont co-financé son développement par des prises de participation directes au capital d'ASML et par des engagements d'achat anticipé, reproduisant un schéma de financement consortial que l'industrie des semi-conducteurs avait inauguré avec SEMATECH (consortium public-privé américain créé en 1987) et que l'institut IMEC (Louvain, fondé en 1984) a prolongé en Europe. Ce mode de financement, dans lequel des concurrents sur le marché aval co-investissent dans leur fournisseur commun en amont, n'a pas d'équivalent dans les autres GPT documentées par ce chapitre~; il constitue une forme de coordination par la \emph{roadmap} qui se distingue tant de la coordination par le marché que de la coordination par l'État.

\dimmark{D9}{*}%
La division géographique du travail qui résulte de cette trajectoire est elle-même un fait historique sans précédent dans l'industrie informationnelle. Le design des puces les plus avancées est concentré aux États-Unis (Apple, Nvidia, Qualcomm, AMD)~; la fabrication est dominée par TSMC à Taïwan et Samsung en Corée du Sud~; l'équipement critique de lithographie est produit aux Pays-Bas (ASML) avec des optiques allemandes (Carl Zeiss SMT) et des lasers allemands (Trumpf). Aucun pays ne maîtrise la chaîne complète. Cette configuration n'est pas une nécessité technique~; c'est le résultat contingent de quarante ans de décisions d'investissement, de politiques industrielles et de spécialisations cumulatives. Elle structure la géopolitique contemporaine des semi-conducteurs~: les restrictions américaines à l'exportation de machines EUV vers la Chine (à partir de 2019 pour les machines les plus avancées), le \emph{CHIPS and Science Act} américain (2022), le \emph{European Chips Act} (2023) et le plan \emph{Made in China 2025} sont autant de réponses à une vulnérabilité que cette division du travail a créée. Le parallèle avec la domination britannique sur la gutta-percha et les câbles sous-marins au XIXe~siècle (§\ref{subsec:materialite}) est structurel, non métaphorique~: dans les deux cas, le contrôle d'un input matériel critique confère un pouvoir de veto sur l'infrastructure informationnelle mondiale. Mais les conditions de production du monopole diffèrent~: le monopole britannique reposait sur l'accès géographique à une matière première extractive~; le monopole d'ASML repose sur l'accumulation de compétences et de coûts fixes irrécouvrables dans un processus de fabrication d'une complexité sans équivalent.

\dimmark{0-TER}{*}%
La conséquence énergétique de cette trajectoire est la plus directement pertinente pour l'argument de la thèse. Chaque machine EUV en fonctionnement consomme de l'ordre d'un mégawatt, soit environ dix fois la consommation d'une machine DUV de génération précédente. Une usine de fabrication de semi-conducteurs avancés (un \emph{fab} aux n{\oe}uds 5~nm ou 3~nm) opère plusieurs dizaines de ces machines simultanément, vingt-quatre heures sur vingt-quatre. La consommation électrique de TSMC représentait en 2023 environ six pour cent de la consommation totale d'électricité de Taïwan, une proportion projetée à 12,5~pour cent en 2025\footnote{Estimation citée par Bloomberg en août 2022 et reprise dans la littérature de suivi sectoriel. La proportion exacte dépend de la trajectoire de déploiement des n{\oe}uds avancés et de l'évolution de la demande électrique totale de Taïwan.}. La poursuite de la miniaturisation, loin de réduire l'empreinte matérielle de l'ICT, déplace l'intensification du produit vers le procédé~: le coût énergétique unitaire de chaque transistor gravé continue de baisser, mais le coût énergétique du système de production qui le rend possible augmente à chaque génération technologique. C'est une externalisation de la contrainte matérielle, du même type que celle que le chapitre a documentée pour le passage du cuivre à la fibre optique (§\ref{subsec:www})~: la contrainte ne disparaît pas~; elle se déplace vers un niveau d'infrastructure moins visible.

\emergence{La trajectoire de la photolithographie révèle que la loi de Moore n'est pas une loi naturelle de dématérialisation~; c'est un dispositif de coordination industrielle dont la prolongation exige une intensification croissante du substrat énergétique et matériel de production. Le monopole d'ASML, la division géographique du travail et la consommation électrique des fabs avancés sont les résultats historiques de cette trajectoire. Pour la thèse, le cas établit que l'intensification énergétique du système informationnel ne se joue pas seulement côté usage (centres de données, terminaux, réseaux) mais aussi côté production (fabrication des puces elles-mêmes), et que les deux dynamiques se composent.}

\footnote{Pour le récit géopolitique de la concentration du secteur, voir \textcite{miller_chip_2022}. Les données techniques sur les machines EUV proviennent des rapports annuels ASML (2020--2024) et des présentations investisseurs. Sur le financement consortial SEMATECH et ses effets sur la structure de l'industrie américaine, voir \textcite{irwin_learning_1994} et \textcite{flamm_targeting_1987}.}


\subsection{Le stockage numérique et la mémoire permanente}\label{subsec:store}

\dimmark{0-TER}{*}%
La section~\ref{subsec:phonographe} avait établi que le régime pré-numérique du stockage est passif~: la cire, le disque, le film fixent l'information sans énergie, et cette information peut être relue des décennies plus tard sans qu'un seul watt ait été consommé entre les deux actes. Ce régime prend fin avec la mémoire à accès dynamique~: la DRAM (Intel~1103, 1970) est la première mémoire dont les cellules doivent être rafraîchies plusieurs centaines de fois par seconde pour maintenir leur contenu~; sans alimentation continue, l'information disparaît. La DRAM n'est pas une amélioration du stockage passif~: c'est une rupture de régime, un changement dans la nature du couplage entre information et énergie. La première s'applique dès lors à la mémoire de travail des ordinateurs. L'étape suivante est la généralisation de ce régime au stockage persistant lui-même.

\dimmark{0-TER}{*}%
Le Web (Berners-Lee, 1991) est l'événement qui accomplit cette généralisation. Un document stocké sur un serveur web n'est pas seulement un fichier~: c'est un fichier auquel n'importe qui peut accéder à n'importe quel moment, depuis n'importe où, sans préavis. Pour que cet accès soit possible, le serveur doit être allumé en permanence. La consommation d'énergie du serveur ne dépend pas du nombre de lectures~: un fichier consulté une fois par an et un fichier consulté un million de fois par jour consomment la même énergie sur le serveur. Ce qui a changé n'est pas la nature de l'information stockée~: c'est l'exigence de disponibilité permanente qui lui est associée, et cette exigence impose un flux d'énergie continu indépendamment de l'usage. C'est précisément le mécanisme identifié dans la section~\ref{sec:convergence}~: STORE couplé à TRANSMIT sur infrastructure \emph{always-on} produit une consommation continue, là où tous les substrats précédents (papier, carton, vinyle, film) ne consommaient qu'à la fabrication et à la lecture.

\dimmark{0-TER}{*}%
\textcite{hilbert_worlds_2011} mesurent la montée en puissance de ce régime~: en 2002, pour la première fois, la capacité de stockage numérique mondiale dépasse la capacité analogique~; en 2007, quatre-vingt-quatorze pour cent de l'information mondiale est numérique en termes de capacité de stockage, et quatre-vingt-dix-neuf virgule neuf pour cent des télécommunications sont numérisées\footnote{Le chiffre de 94~\% concerne le stockage~; le 99,9~\% concerne les télécommunications. Ces deux grandeurs ne doivent pas être confondues. L'étude de \textcite{hilbert_worlds_2011} couvre la période 1986--2007 et ne fournit pas de données post-2007.}. Le coût énergétique de ce basculement est invisible dans leurs données~: Hilbert et López mesurent les capacités fonctionnelles (bits par seconde de transmission, bits de stockage, MIPS de calcul) mais pas les kilowattheures associés. L'angle mort est structurel~: leur cadre analytique, comme celui de Shannon (1948) dont il procède, traite l'information indépendamment de son substrat énergétique.


\subsection{L'intelligence artificielle générative et les lois d'échelle}\label{subsec:genai}

\dimmark{0-TER}{*}%
Le basculement de régime de la couche compute s'opère avec la diffusion des lois d'échelle (\emph{scaling laws}) dans l'entraînement des grands modèles de langage. \textcite{kaplan_scaling_2020} établissent empiriquement que la performance prédictive d'un modèle de langage évolue comme une loi de puissance du compute investi, du nombre de paramètres et de la taille du corpus d'entraînement. \textcite{hoffmann_training_2022} précisent les rapports optimaux entre ces trois variables (les «~lois Chinchilla~\fg). Ce que ces résultats impliquent pour l'économie du système informatique est considérable~: améliorer la frontière des performances nécessite structurellement \emph{plus} de compute, pas moins. Il n'existe pas de substitut organisationnel, algorithmique ou architectural qui permettrait d'atteindre le même niveau de performance sans augmenter le compute~--- du moins dans le paradigme Transformer qui domine l'industrie depuis 2017. Le rapport improvement/énergie, que tous les cas précédents avaient vu évoluer dans le même sens (plus de performance pour la même énergie ou moins), s'inverse pour cette couche spécifique.

\dimmark{0-TER}{*}%
La traduction en puissance de rack est documentable. Les serveurs équipés de GPU Ampere (2020) consommaient environ treize kilowatts par rack~; les serveurs Blackwell (2024) en consomment environ cent trente~; les feuilles de route prévoient un mégawatt par rack pour les générations Rubin (2026 et au-delà)\footnote{Ces chiffres sont issus des rapports techniques et des annonces constructeurs (NVIDIA). \textcite{sun_chip_2024} et \textcite{roussilhe_phase_2025} les documentent dans le contexte de la rupture de la loi de Dennard.}. La traduction en coût énergétique par usage est tout aussi significative~: une requête ChatGPT exige environ dix fois plus d'énergie qu'une recherche Google de génération précédente. Pour la première fois dans l'histoire traversée par ce chapitre, l'énergie est un coût marginal de production du service informationnel~: chaque inférence a un coût électrique directement proportionnel à sa longueur et à la taille du modèle qui la produit.

\dimmark{D2/D8}{+}%
La stratification empirique crée une asymétrie analytique essentielle. Les couches réseau et stockage classique restent en régime Koomey~: \textcite{mytton_network_2024} documentent que la consommation énergétique des réseaux n'est pas proportionnelle au volume de données transmises~; chez Telefónica, le trafic a augmenté de cinq cent soixante-cinq pour cent entre 2015 et 2021 tandis que la consommation baissait de sept virgule deux pour cent. Le coût énergétique par bit transmis continue de décroître d'environ trente pour cent par an. Il n'y a pas de rupture dans ces couches. La rupture est localisée dans le compute IA-intensif~--- et c'est cette localisation qui la rend identifiable, mesurable, et modélisable dans le cadre IMACLIM. Parce que la couche qui a basculé est celle qui concentre l'essentiel de la croissance de l'investissement~--- plus de sept cent quinze milliards de dollars de dépenses d'investissement annoncées par les grands opérateurs pour 2026~--- elle tire l'ensemble du système malgré sa localisation dans un segment technique spécifique\footnote{Pour une analogie structurelle~: dans les années 1900, l'automobile représentait un segment étroit de l'économie de la mobilité, dominée par le cheval et le rail. Parce qu'elle concentrait l'investissement et l'innovation, elle a reconfiguré l'ensemble du système de transport, de l'urbanisme et de l'infrastructure pétrolière. Le compute IA occupe une position structurellement analogue.}.

\dimmark{D8}{*}%
La traduction la plus spectaculaire de cette rupture localisée est l'émergence des contrats d'approvisionnement en énergie nucléaire. Microsoft signe un accord d'achat d'électricité sur vingt ans avec Constellation Energy pour le redémarrage de la centrale Three Mile Island (835~MW)~; Amazon contractualise 1,92~GW depuis la centrale de Susquehanna sur dix-sept ans~; Meta signe un accord de 1,100~MW depuis la centrale Clinton. Plus de dix gigawatts de capacité nucléaire ont été contractés par les grandes firmes numériques en 2024--2025 aux seuls États-Unis \parencite{iea_energy_ai_2025}. Pour la première fois dans l'histoire de l'industrie tech, des entreprises dont le métier est le traitement de l'information signent des contrats de \emph{production} d'énergie sur quinze à vingt ans~: les horizons d'investissement des deux secteurs se synchronisent, et le couplage financier-temporel entre compute et énergie acquiert une matérialité contractuelle qu'aucune phase précédente de l'histoire informationnelle n'avait produite.


\subsection{La blockchain et les cryptomonnaies}\label{subsec:blockchain}

\dimmark{0-TER}{*}%
Bitcoin (Nakamoto, 2009) est le cas extrême du couplage calcul-énergie~: son protocole de preuve de travail (\emph{proof-of-work}) exige que la sécurité du réseau soit achetée en kilowattheures. Le mécanisme est délibéré~: pour qu'une transaction soit validée, des machines doivent résoudre un problème cryptographique dont la difficulté est ajustée de façon à ce que la résolution dure environ dix minutes, quelle que soit la puissance de calcul disponible sur le réseau. Augmenter la puissance de calcul ne réduit pas la consommation d'énergie par transaction~: cela augmente la difficulté et maintient la consommation énergétique totale. L'indice de Cambridge sur la consommation électrique du Bitcoin (\emph{Cambridge Bitcoin Electricity Consumption Index}) estime la consommation annuelle du réseau Bitcoin à cent cinquante à deux cents~TWh~--- un ordre de grandeur comparable à la consommation d'un pays européen de taille moyenne. Ethereum a migré vers un protocole de preuve d'enjeu (\emph{proof-of-stake}) en 2022, réduisant sa consommation de plus de quatre-vingt-dix-neuf pour cent~: ce contraste entre les deux protocoles illustre que le couplage compute-énergie n'est pas une nécessité physique pour toute blockchain, mais un choix d'architecture de gouvernance.

\dimmark{D10/0-TER}{+}%
Le cas mérite d'être pris pour ce qu'il est. La preuve de travail ne consomme pas de l'électricité par accident d'ingénierie~; elle transforme délibérément une dépense computationnelle vérifiable en mécanisme de sécurité, c'est-à-dire qu'elle remplace une autorité centrale par un coût physique imposé aux participants. La dépense d'énergie n'est pas un résidu du protocole~: elle est la condition par laquelle le protocole rend coûteuse l'attaque, et c'est ce qui distingue le registre distribué d'un simple fichier répliqué. Mais cette propriété fait de Bitcoin un bon révélateur et un mauvais modèle général. Elle révèle qu'une architecture informationnelle peut convertir de l'énergie en confiance distribuée, c'est-à-dire acheter de la rareté dans un monde numérique où la copie est presque gratuite~; elle ne dit pas que tout système informationnel doit s'organiser ainsi. Le cas doit rester dans le chapitre, mais comme cas limite~: il montre une modalité du couplage information-énergie inscrit dans la règle de gouvernance elle-même, non la loi générale du numérique contemporain.


\subsection{L'IA agentique et l'automatisation cognitive}\label{subsec:agentique}

\dimmark{D5/D10}{+}%
La génération de texte, d'image ou de code ne clôt pas la trajectoire de l'intelligence artificielle contemporaine~; elle en déplace le centre de gravité. Tant que le modèle produit une réponse à une requête, le calcul reste attaché à un acte humain discret~: l'utilisateur demande, la machine infère, le résultat revient dans l'interface. Avec ce qu'on appelle l'IA agentique, le calcul change de statut. Il ne sert plus seulement à produire une sortie informationnelle~; il sert à enchaîner des opérations, à choisir des outils, à interroger des bases, à écrire du code, à vérifier partiellement ses propres résultats, et à relancer le processus tant que l'objectif n'est pas atteint. Le modèle devient alors moins un générateur qu'un opérateur.

Ce déplacement ne doit pas être décrit comme une autonomie au sens fort. Les agents contemporains restent fragiles, dépendants de leur environnement logiciel, exposés aux erreurs d'exécution, aux hallucinations et aux objectifs mal spécifiés\footnote{La littérature récente sur les agents fondés sur les grands modèles de langage reste instable~: elle documente des progrès dans l'usage d'outils, la planification et la décomposition de tâches, mais aussi des taux d'échec élevés dès que les environnements deviennent longs, ouverts ou faiblement spécifiés. Le point retenu ici n'est donc pas l'autonomie des agents, qui demeure une promesse technique, mais la modification de l'unité économique pertinente~: une tâche humaine peut être transformée en séquence d'appels de modèles, d'outils et de vérifications.}. Mais la rupture n'a pas besoin d'être parfaite pour avoir une portée économique~: il suffit que l'agent transforme une requête humaine en trajet computationnel composé de plusieurs inférences enchaînées pour que le rapport entre acte humain et budget de calcul cesse d'être de un pour un. Le rapprochement avec les techniques antérieures de l'information~--- le télégraphe qui n'avait pas supprimé tous les délais mais certains délais pertinents, la tabulatrice qui n'avait pas supprimé le travail administratif mais déplacé son goulot d'étranglement~--- ne porte pas la preuve~; il sert seulement à indiquer que la structure du déplacement n'est pas neuve.

\dimmark{D5/0-TER}{+}%
C'est ici que le changement devient analytique. Le numérique de plateforme restait borné par le temps humain~: on regardait une vidéo, on écrivait un message, on cherchait une information. Le numérique agentique relâche partiellement cette borne, car une tâche humaine peut devenir un processus computationnel composé d'opérations intermédiaires. La contrainte ne disparaît pas~; elle change de lieu. Elle passe du temps de l'utilisateur au budget de calcul, de l'attention humaine à la capacité d'orchestration, de la rareté cognitive à la rareté infrastructurelle. En termes contemporains, on dirait que la demande devient partiellement \emph{machine-to-machine}~; cette formulation reste conjecturale tant que les mesures empiriques manquent. Le chapitre~\ref{ch:structures} aura à déterminer si cette modification de l'unité économique pertinente constitue un effet de rebond d'une nature inédite (D5), ou simplement une croissance ordinaire de l'usage telle que la littérature contemporaine sur le rebond la modélise déjà.


\subsection{L'Internet des objets et l'\emph{edge computing}}\label{subsec:iot}

\dimmark{D9/0-TER}{+}%
Le déploiement des objets connectés à partir des années 2010~--- compteurs intelligents, capteurs industriels, dispositifs domestiques, terminaux logistiques~--- prolonge un mouvement déjà engagé par la téléphonie mobile~: l'informatique sort des centres de calcul pour se diffuser dans l'environnement matériel ordinaire. Le décompte précis du parc d'équipements reste un objet de communication industrielle plus que de mesure scientifique~; les rapports des équipementiers réseau (Cisco, Ericsson) avancent des chiffres de plusieurs dizaines de milliards d'objets connectés à la fin de la décennie 2020, mais leur fiabilité est limitée par la définition retenue du périmètre et par l'absence de méthodologie partagée\footnote{Les rapports industriels (Cisco \emph{Annual Internet Report}, Ericsson \emph{Mobility Report}, ITU \emph{Measuring the Information Society}) doivent être traités comme des documents de marché plus que comme des sources de mesure. Leur fonction est de coordonner les anticipations d'investissement du secteur, et les ordres de grandeur qu'ils annoncent doivent être tenus comme des indicateurs de tendance, non comme des résultats consolidés.}. Ce qui importe pour le présent chapitre n'est pas le nombre exact de capteurs mais la transformation qualitative qu'ils induisent~: le calcul migre vers la périphérie du réseau parce que la latence, la bande passante et la confidentialité ne permettent plus toujours d'acheminer chaque donnée jusqu'à un centre de données distant.

\dimmark{0-TER}{+}%
La promesse énergétique de cette migration vers l'\emph{edge} est ambiguë. Réduire le volume de données transmises diminue la consommation des réseaux longue distance~; calculer localement consomme moins par opération que d'attendre l'aller-retour vers le cloud~; chiffrer et préfiltrer à la source économise du stockage en aval. Mais la multiplication des nœuds de calcul périphériques crée elle-même une consommation distribuée, dont l'agrégation est moins visible que celle des centres de données mais n'est pas pour autant négligeable~: chaque équipement maintient une alimentation, une mémoire, une connectivité, et une part de calcul. Surtout, les usages les plus dépendants de l'IoT massif~--- véhicule autonome, ville intelligente, jumeau industriel~--- réagrègent les flux périphériques dans des modèles entraînés et exécutés en cloud, où la rupture identifiée pour la couche IA générative s'applique de plein droit. La consommation totale n'est donc pas réduite par l'\emph{edge}~; elle est redistribuée géographiquement entre le terminal, l'antenne, le nœud de bord et le centre de données. Le Ch.~\ref{ch:biophysique} aura à mesurer cette redistribution.


\subsection{L'IA physique, la robotique et les véhicules autonomes}\label{subsec:physical_ai}

\dimmark{0-TER}{+}%
Une frontière supplémentaire se déplace lorsque le calcul sort de l'écran et entre dans l'action physique. Les véhicules autonomes (Waymo en exploitation commerciale limitée à San Francisco et Phoenix, Tesla Full Self-Driving en bêta sur flotte client), les robots humanoïdes (Tesla Optimus, Figure, Apptronik), les robots mobiles industriels (Boston Dynamics, Agility Robotics) et les systèmes d'assistance physique constituent une famille hétérogène mais structurellement comparable~: ce sont des plateformes informationnelles dont la sortie n'est plus un fichier ou un signal, mais un mouvement dans un environnement matériel. Cette caractéristique modifie le profil de consommation. La machine embarque un budget de calcul significatif pour percevoir, planifier et agir en temps quasi réel~; elle dépend en parallèle d'un entraînement en cloud qui ingère des volumes massifs de données capturées sur le terrain.

\dimmark{0-TER}{+}%
Le statut de ce résultat doit rester prudent. La plupart de ces plateformes sont en déploiement limité ou en démonstration~: les promesses industrielles sont nombreuses, les résultats publiés sont fragmentés, et l'écart entre prototype et industrialisation est documenté de longue date dans la robotique. Ce que l'on peut tenir, en revanche, est la structure du couplage énergétique. L'IA physique combine trois sources de consommation~: l'inférence embarquée, dont la puissance dépend du compromis entre autonomie d'action et coût du composant~; l'entraînement en cloud, qui relève du régime décrit à la section~\ref{subsec:genai}~; et l'entraînement en simulation, où des moteurs physiques distribués (NVIDIA Omniverse et plateformes équivalentes) génèrent des environnements virtuels destinés à compléter les données réelles. Cette troisième source est peu mesurée mais non négligeable~: simuler le monde physique pour entraîner des machines physiques transfère une partie de l'apprentissage de l'extraction de données vers la génération de scénarios, déplaçant ainsi le coût énergétique sans nécessairement le réduire.


\subsection{Les jumeaux numériques et la simulation à grande échelle}\label{subsec:digital_twins}

\dimmark{0-TER}{+}%
Le jumeau numérique désigne une copie informationnelle d'un système physique, tenue synchronisée avec lui par un flux continu de données issues de capteurs et par un modèle qui en simule le comportement. Le concept a été formulé au début des années 2000 dans le domaine de la gestion du cycle de vie produit, puis étendu à des objets plus vastes~: usines (Siemens Xcelerator), réseaux électriques, infrastructures urbaines, systèmes climatiques (NVIDIA Earth-2, programme \emph{Destination Earth} de l'Union européenne). Son intérêt analytique pour le présent chapitre tient moins à la technologie elle-même qu'à son régime de fonctionnement~: à la différence d'un modèle numérique classique, qu'on exécute ponctuellement pour répondre à une question, le jumeau numérique est censé tourner en continu pour maintenir sa correspondance avec le réel.

\dimmark{0-TER}{+}%
Cette continuité a une conséquence énergétique de principe~: un jumeau opérationnel consomme du calcul même lorsque le système physique qu'il représente n'évolue pas, parce que la synchronisation impose un cycle permanent d'ingestion, de simulation et de vérification. Pour les jumeaux d'objets industriels limités, l'ordre de grandeur reste maîtrisable~; pour les jumeaux climatiques ou terrestres à haute résolution, il devient un sujet en soi, puisque la simulation d'un système physique à grande échelle pour le suivre en temps quasi réel mobilise des infrastructures de calcul de classe supercalculateur. Le déploiement de ces objets est encore largement industriel et expérimental, et les mesures consolidées de leur empreinte sont rares~; le présent chapitre ne peut donc qu'enregistrer la trajectoire, pas la quantifier. Le Ch.~\ref{ch:biophysique} aura à déterminer si les jumeaux numériques constituent une catégorie d'usage suffisamment volumineuse pour modifier la trajectoire agrégée de consommation, ou s'ils restent un cas particulier dont la portée est plus symbolique qu'énergétique.


\subsection{Les réseaux 5G/6G et l'infrastructure de transmission}\label{subsec:5g}

\dimmark{D9/0-TER}{+}%
La cinquième génération de réseaux mobiles, déployée à partir de 2019, se distingue de la quatrième par une densification des stations de base, l'usage de bandes de fréquence millimétriques (entre 24 et 39~GHz) et une intégration plus poussée des fonctions de calcul dans le réseau lui-même. Le coût énergétique de ce déploiement est souvent présenté comme une rupture~: une station de base 5G consomme deux à trois fois plus qu'une station 4G équivalente. La formulation est exacte mais incomplète, et c'est précisément l'incomplétude qui doit être tenue~: la consommation par station ne dit rien de la consommation par bit transmis. Les bandes millimétriques sont plus directionnelles et requièrent des cellules plus petites, mais elles transportent aussi un volume de données beaucoup plus important par unité de spectre. La conséquence agrégée n'est donc pas mécaniquement une augmentation de la consommation~: elle dépend du rapport entre la croissance des stations et la croissance des données transportées.

\dimmark{0-TER}{+}%
C'est exactement le mécanisme que \textcite{mytton_network_2024} ont identifié pour les générations précédentes~: la consommation énergétique des réseaux n'est pas proportionnelle au trafic, parce que les améliorations d'efficacité par bit absorbent une part substantielle de la croissance des usages. Tant que ce mécanisme tient, les réseaux mobiles continuent d'opérer dans le régime Koomey décrit à la section~\ref{subsec:genai}, et leur empreinte énergétique reste un sujet de gestion plutôt qu'un sujet de rupture. Les questions ouvertes pour le Ch.~\ref{ch:biophysique} portent sur les configurations qui pourraient le faire céder~: la densification radicale anticipée pour la 6G, l'intégration calcul-réseau qui ferait remonter une part du compute dans l'infrastructure de transmission, ou le déploiement parallèle de constellations satellitaires en orbite basse (Starlink, Kuiper) dont le coût énergétique par bit reste mal documenté et qui ajoutent une couche d'infrastructure dont la matérialité~--- lancements, durée de vie courte, remplacement permanent~--- est d'un autre ordre que celle des antennes terrestres.


\subsection{La fenêtre 2008 et les trois cadres régulatoires post-crise}\label{subsec:fenetre_2008}

\dimmark{D8/0-GOV}{+}%
La crise financière de 2008 mérite d'être nommée dans l'histoire de la trajectoire numérique non pour ses causes proximales, qui relèvent de l'analyse macro-financière, mais pour ce qu'elle a ouvert et refermé à propos du couplage ICT-économie. La crise survient au moment exact où la phase \emph{synergy} commence à se stabiliser (§\ref{subsec:turning_point})~: le \emph{cloud computing} démarre commercialement, le modèle \emph{surplus comportemental} se généralise, l'iPhone et l'App Store sont lancés. Plusieurs propositions de politique publique formulées à ce moment auraient pu réorienter cette stabilisation. \textcite{hines_green_2008} et la \emph{Green New Deal Group} britannique proposent en juillet 2008 un \emph{Green New Deal} qui associe décarbonation rapide de l'économie, régulation financière renforcée et investissement public massif dans les infrastructures bas-carbone. La proposition ne sera retenue ni par les administrations américaine et britannique post-crise, ni par les institutions européennes, dont la réponse principale combine plans de relance limités, politique monétaire ultra-accommodante (\emph{quantitative easing} à partir de 2009) et maintien de l'autorégulation industrielle dans le numérique. \textcite{varoufakis_techno_2023} a proposé du résultat agrégé l'expression de \emph{techno-féodalisme}~: la combinaison de \emph{cheap credit} (taux quasi nuls pendant une décennie) et de captation rentale par les plateformes numériques produit une expansion sans précédent des centres de données et des firmes \emph{big tech}, dont la capitalisation atteint en 2020 un niveau que la «~révolution dot-com~» de 1999 n'avait jamais approché. Le lien causal direct entre crise 2008 et expansion data centers reste à quantifier~; ce que documente l'historiographie disponible est que la fenêtre d'une bifurcation régulatoire et infrastructurelle s'est ouverte en 2008--2009 et s'est refermée sans qu'aucune des options alternatives portées à ce moment ne soit retenue~\parencite{lundvall_techno_2017}.

\dimmark{0-GOV/D10}{+}%
C'est dans le cadre laissé par cette fermeture que se constituent, entre 2010 et 2024, les trois grands cadres régulatoires du numérique contemporain. \textcite{bradford_digital_2023} les caractérise par leur ancrage normatif différencié~: le modèle américain reste \emph{market-driven}, prolongeant l'autorégulation issue du \emph{Framework} de 1997 (§\ref{subsec:www})~; le modèle européen adopte une orientation \emph{rights-driven} dont les jalons sont le Règlement général sur la protection des données (RGPD, 2018), le \emph{Digital Services Act} et le \emph{Digital Markets Act} (2022), et l'\emph{AI Act} (2024)~; le modèle chinois est \emph{state-led} et combine contrôle politique du contenu, soutien stratégique à l'industrie nationale et régulation antimonopolistique à géométrie variable. \textcite{moldicz_techno_2022} documente la course à la \emph{souveraineté technologique} qui se déploie entre ces trois pôles après 2018~: contrôles américains à l'exportation des semi-conducteurs avancés, plan \emph{Made in China 2025}, plans européens \emph{Chips Act} et \emph{Cloud sovérain}. Ces trois cadres ne convergent pas vers une régulation mondiale unifiée du couplage ICT-énergie~; ils convergent vers une fragmentation régulatoire qui rend plus difficile, plutôt que plus facile, la formulation de politiques visant le couplage lui-même.


\subsection{La critique des \emph{twin transitions} et les niches alternatives}\label{subsec:critique_twin}

\dimmark{0-TER/D10}{+}%
Le discours dominant de la décennie 2020 sur la digitalisation est celui des \emph{twin transitions}~: l'idée portée par la Commission européenne et par les organisations internationales (OCDE, AIE) que la transition numérique et la transition écologique sont structurellement complémentaires et se renforcent mutuellement~\parencite{european_commission_twin_2022}. La littérature académique récente est venue questionner cette hypothèse de complémentarité. \textcite{santarius_digital_2022} ont proposé le concept de \emph{digital sufficiency} pour interroger l'hypothèse que la digitalisation contribuerait nécessairement à la décarbonation. \textcite{creutzig_digitalization_2022} dans la \emph{Annual Review of Environment and Resources} produisent un état de l'art qui documente la coexistence d'effets décarbonants (gestion fine de l'énergie, optimisation logistique) et d'effets rebond significatifs (croissance des usages absorbant les gains d'efficacité). \textcite{bennich_untangling_2025} et \textcite{mouthaan_systemic_2023} démontrent par revue systématique que l'hypothèse de complémentarité mécanique n'est pas étayée par les données disponibles, et que les conditions sous lesquelles digitalisation et décarbonation coopèrent effectivement relèvent de configurations politiques précises, non d'une dynamique structurelle automatique. \textcite{makitie_digital_2023} prolongent en montrant que la conception même d'une \emph{digital green transition} comme objet politique unifié pose des problèmes analytiques mal résolus. Le chapitre~\ref{ch:biophysique} aura à reprendre ce dossier en évaluant si, sous quelles conditions et dans quelles proportions, le déploiement numérique observé depuis 2007 peut effectivement contribuer à la décarbonation des autres secteurs ou, au contraire, pèse sur les budgets carbone disponibles.

\dimmark{D10}{+}%
La critique des \emph{twin transitions} ne dit pas qu'aucune alternative au déploiement numérique dominant n'existe~; elle dit que les alternatives existantes restent confinées à des niches dont la portée systémique reste à établir. Plusieurs configurations méritent d'être nommées, sans qu'on en surestime la portée ni qu'on les présente comme un programme de substitution. Le projet \emph{Fairphone} (lancé en 2013 aux Pays-Bas) propose un téléphone modulaire conçu pour la réparation et l'allongement de la durée d'usage~; \textcite{haucke_smartphone_2018} documente sa trajectoire économique limitée mais analytiquement instructive. Le projet \emph{Te Hiku Media} en Nouvelle-Zélande développe depuis 2018 des modèles de reconnaissance vocale en langue maori sous gouvernance communautaire stricte, et le collectif \emph{DAIR} (\emph{Distributed AI Research Institute}) fondé par Timnit Gebru porte une alternative au modèle \emph{scale-at-all-costs} de l'IA contemporaine~\parencite{hao_empire_2025}. La municipalité de Barcelone a développé entre 2015 et 2023 le projet \emph{DECODE} de souveraineté numérique citoyenne~\parencite{smith_smart_2021}, et la plateforme \emph{Energy Community Platform} de l'Union européenne (2024) soutient le déploiement de coopératives énergétiques utilisant des outils numériques sous gouvernance partagée~\parencite{energy_community_2024}. Ces configurations partagent un trait analytique~: elles découplent l'usage du numérique du couplage \emph{surplus comportemental + scaling + concentration} qui structure la trajectoire dominante depuis 2007. Le chapitre~\ref{ch:biophysique} et le chapitre~\ref{ch:experimentation} auront à déterminer si ces niches peuvent constituer des configurations de référence pour l'analyse de trajectoires alternatives, ou si leur marginalité même est un résultat de structure qu'aucune dynamique endogène du système actuel ne paraît susceptible de modifier.


\subsection{La limite thermodynamique~: le principe de Landauer}\label{subsec:landauer}

\dimmark{0-TER}{+}%
L'histoire que ce chapitre a suivie pourrait donner l'impression que l'information échappe toujours à ses supports. Un code remplace un signal optique~; un fil remplace un messager~; une carte perforée remplace une colonne de chiffres~; un transistor remplace une lampe~; un modèle statistique remplace une règle explicite. À chaque étape, le support change, et l'on pourrait croire que la trajectoire conduit vers une abstraction croissante. Le principe de Landauer rappelle que cette abstraction a un plancher.

\textcite{landauer_irreversibility_1961} établit que l'effacement irréversible d'un bit d'information entraîne une dissipation minimale d'énergie proportionnelle à $kT\ln 2$. La formule est minuscule à l'échelle d'un bit~; elle est décisive à l'échelle du principe. Elle ne dit pas que les ordinateurs contemporains sont proches de cette limite, ni que l'énergie des centres de données serait réductible à l'effacement logique. Elle dit quelque chose de plus profond pour l'histoire économique de l'information~: le calcul n'est pas seulement une opération formelle, car toute opération physiquement réalisée modifie un support matériel et dissipe de l'énergie lorsqu'elle est irréversible. \textcite{bennett_logical_1973,bennett_thermodynamics_1982} a ensuite montré qu'un calcul logiquement réversible peut contourner cette dissipation minimale, mais au prix d'une autre contrainte~: conserver l'information au lieu de l'effacer, c'est déplacer la charge vers la mémoire, la complexité et la gestion des états.

\dimmark{0-TER}{+}%
Ce détour thermodynamique ne doit pas être transformé en causalité historique directe\footnote{La littérature sur Landauer est traversée par une difficulté terminologique~: l'entropie informationnelle de Shannon, l'entropie thermodynamique et l'entropie algorithmique ne désignent pas la même grandeur, même si les analogies entre elles ont été fécondes. Bennett a précisément contribué à clarifier le paradoxe du démon de Maxwell en distinguant l'acquisition d'information de l'effacement de mémoire. Le chapitre n'utilise donc pas Landauer pour expliquer la consommation actuelle des centres de données~; il l'utilise pour borner conceptuellement l'idée d'un calcul immatériel.}. Landauer n'explique ni la loi de Moore, ni l'économie des semi-conducteurs, ni la croissance contemporaine des infrastructures de calcul. Il fournit une borne conceptuelle. Or cette borne suffit à rompre avec l'un des gestes fondateurs de la théorie de l'information~: la séparation entre information et énergie. Shannon avait raison de séparer le contenu sémantique du signal pour construire une théorie mathématique de la communication~; il n'en résultait pas que l'information réalisée puisse être séparée indéfiniment de son support. La fécondité de la séparation formelle a produit, par retour, une occultation matérielle.

Cette conclusion ne signifie pas que toute augmentation de calcul conduit mécaniquement à une augmentation proportionnelle d'énergie. L'histoire récente de l'informatique est aussi celle d'améliorations d'efficacité, d'architectures spécialisées et de gains matériels qui déplacent sans cesse le rapport entre opération logique et dissipation effective. Mais elle interdit de traiter le calcul comme une activité purement immatérielle. Le bit n'a pas de joule dans la théorie de Shannon~; le bit effacé en a un dans le monde de Landauer. Le Ch.~\ref{ch:structures} pourra alors reprendre la question sous une autre forme~: si le calcul possède un plancher physique, quelles structures économiques déterminent la vitesse à laquelle les sociétés s'en approchent, s'en éloignent, ou l'oublient~?


\begin{table}[htbp]
\caption{Dimensions économiques de la période contemporaine (§\ref{sec:contemporain}).}\label{tab:contemporain-dimensions}
\centering\small
\renewcommand{\arraystretch}{1.15}
\begin{tabularx}{\textwidth}{@{}cXc@{}}
\toprule
\textbf{Dim.} & \textbf{Ce que §\ref{sec:contemporain} établit} & \textbf{Niveau} \\
\midrule
D2 & Concentration cloud (AWS, Azure, GCP)~; monopole ASML sur la lithographie EUV~; modèle Insull reproduit à l'échelle mondiale. & ★ \\
D3 & Externalisation du compute vers le cloud~: capital fixe $\rightarrow$ coût variable~; découplage comptable entre valeur logicielle et coût énergétique d'usage. & + \\
D4 & Déplacement du travail vers l'orchestration et la supervision (IA agentique)~; division internationale du travail semi-conducteur (design US, fab Taïwan, équipement NL). & + \\
D5 & Surplus comportemental (\textcite{zuboff_surveillance_2019})~; rebond potentiel de l'IA agentique (tâche $\rightarrow$ séquence d'inférences)~; scaling laws = amélioration $\propto$ compute. & ★ \\
D6 & AdWords/AdSense (2002--2003)~: marché de l'attention opérationnalisé~; requête ChatGPT $\approx$ 10$\times$ énergie recherche Google. & + \\
D8 & Contrats nucléaires 15--20 ans (Microsoft/Three Mile Island, Amazon/Susquehanna, Meta/Clinton)~: horizons d'investissement ICT et énergie synchronisés. Capex 2025 $>$ 320~Mrd\$. & ★ \\
D9 & Division géographique du travail semi-conducteur (US/TW/KR/NL/DE)~; restrictions export EUV vers Chine~; CHIPS Act, European Chips Act, Made in China 2025. & ★ \\
D10 & Loi de Moore comme dispositif de coordination (ITRS/IRDS)~; scaling laws IA comme nouvelle feuille de route~; niches alternatives (Fairphone, Te Hiku, DECODE). & ★ \\
\midrule
0-TER & Rupture Dennard ($\sim$2005)~; datacenters 415~TWh (2024) $\rightarrow$ 945~TWh (2030)~; TDP GPU 700~W (2024) $\rightarrow$ 1\,400~W (2026)~; EUV $\approx$ 1~MW/machine~; TSMC = 6~\% élec. Taïwan~; Bitcoin $\approx$ 150--200~TWh/an~; Landauer = plancher thermodynamique. & ★ \\
0-GOV & Fenêtre 2008 refermée sans bifurcation~; trois cadres régulatoires (\textcite{bradford_digital_2023})~: US \emph{market-driven}, EU \emph{rights-driven}, CN \emph{state-led}~; critique des \emph{twin transitions}. & ★ \\
\bottomrule
\end{tabularx}
\renewcommand{\arraystretch}{1.0}
\end{table}


% =============================================================================
